% UBT-AI-PROVENANCE-BEGIN
% schema: ubt-ai-provenance/v1
% tier: C_working
% ai_assistance: disclosed
% human_review: risk-based
% editorial_responsibility: Ing. David Jaroš
% policy: ../../../AI_PROVENANCE.md
% notice: Working material; exhaustive human review is not claimed.
% UBT-AI-PROVENANCE-END

\chapter{Covariant derivation, connections and pure UBT geometry}
\label{ch:covariant-tetrad}

This chapter explains the geometric core of UBT in the form closest to the
original intuition of the theory:
\[
 \Theta\quad\longrightarrow\quad E_\mu:=D_\mu\Theta
 \quad\longrightarrow\quad
 \{E_\mu,E_\nu\}_\sharp,\ [D_\mu,D_\nu].
\]
The anticommutator carries the metric, while the commutator of covariant
derivatives carries curvature.  No embedding map, compact-fiber average, or
observer projection is needed in the minimal local construction.

\section{Ordinary, partial and covariant derivatives}
For an ordinary scalar function $f(x)$, the derivative $df/dx$ measures how the
number $f$ changes when $x$ changes.  For a field on spacetime we use partial
derivatives,
\[
 \partial_\mu\Theta:=\frac{\partial\Theta}{\partial x^\mu},
 \qquad \mu=0,1,2,3.
\]
A partial derivative compares components written in a fixed basis.  This is
sufficient if the same basis can be used everywhere.  Geometry becomes subtler
when the local frame itself rotates or boosts from point to point.

A simple analogy is a vector written in a rotating coordinate frame.  Even a
physically constant vector acquires changing components because the basis is
changing.  The covariant derivative corrects for this basis change:
\[
 D_\mu\Theta=\partial_\mu\Theta+\rho_*(\Omega_\mu)\Theta.
\]
Here $\Omega_\mu$ is the connection and $\rho_*$ tells us how the connection
acts in the representation of $\Theta$.  For biquaternions, a left action, a
right action, and a two-sided action are genuinely different and must not be
silently interchanged.

\section{What is a connection?}
A connection is not a force by itself.  It is a rule for comparing local frames
at neighboring points.  The symbol $\Omega_\mu$ is the UBT internal or
biquaternionic-frame connection.

Christopher symbols,
\[
 \Gamma^\rho{}_{\mu\nu},
\]
are closely related, but they act on coordinate indices.  For example,
\[
 \nabla_\mu V^\rho=\partial_\mu V^\rho
 +\Gamma^\rho{}_{\mu\sigma}V^\sigma.
\]
Thus:
\begin{itemize}
 \item $\Gamma^\rho{}_{\mu\nu}$ tells us how the coordinate basis changes;
 \item $\omega_\mu{}^a{}_b$, and its spin/biquaternionic lift $\Omega_\mu$,
       tell us how the local Lorentz frame changes.
\end{itemize}
They are related by the tetrad compatibility identity
\[
 \partial_\mu e_\nu{}^a-\Gamma^\rho{}_{\mu\nu}e_\rho{}^a
 +\omega_\mu{}^a{}_b e_\nu{}^b=0.
\]
They are therefore two descriptions of the same geometric transport, not the
same coefficients.  In particular,$\Gamma=\operatorname{Re}\Omega$ is not a
valid identification.

\section{Why the classic connection is not another free field}
For a nondegenerate tetrad, metric compatibility and zero torsion determine a
unique Lorentz connection:
\[
 \boxed{
 \mathring\omega_\mu{}^a{}_b
 =e_b{}^\nu\left(\mathring\Gamma^\rho{}_{\mu\nu}(g)e_\rho{}^a
 -\partial_\mu e_\nu{}^a\right).}
\]
The circle marks the torsion-free Levi--Civita connection.  Thus, in the
ordinary torsion-free GR branch, the frame connection is computed from the
tetrad just as the Christoffel symbols are computed from the metric.  A local
Lorentz rotation changes the numerical coefficients of the tetrad and the
connection but does not change the geometry.

\subsection{Track and counter}
The torsion two-form is
\[
 T^a=de^a+\omega^a{}_b\wedge e^b.
\]
A general metric-compatible connection can be written
\[
 \boxed{\omega=\mathring\omega(e)+K(T),}
\]
where $K$ is the contorsion.  With
$T^a=\tfrac12T^a{}_{bc}e^b\wedge e^c$ and the convention above,
\[
 \boxed{
 K_{abc}=\frac12\left(T_{cab}-T_{abc}-T_{bca}\right).}
\]
This formula has an important conceptual consequence: once the tetrad and the
torsion are specified, the metric-compatible connection is unique.  The open
full-UBT question is therefore not an arbitrary leftover $\Omega_\mu$; it is
the dynamical law selecting $T$.

If the UBT vacuum equations imply $T=0$, the connection reduces to the
Levi--Civita spin connection.  If spin or another biquaternionic source creates
torsion, the same theorem reconstructs the connection from that torsion.

\section{Flat space and explicit UBT representative}
In flat Minkowski spacetime, using inertial Cartesian coordinates and a constant
local frame, one can choose
\[
 \Gamma^\rho{}_{\mu\nu}=0,
 \qquad \Omega_\mu=0,
 \qquad D_\mu=\partial_\mu.
\]
This statement can be strengthened from a choice of gauge to an explicit
single-field represent.  Easy
\[
 \boxed{
 \Theta_{\mathrm{SR}}(x)=\Theta_0+\sqrt{\mathcal N_0}
 \left(i x^0\mathbf1+x^k\mathbf e_k\right).}
\]
Then
\[
 \mathcal N_0^{-1/2}\partial_0\Theta_{\mathrm{SR}}=i\mathbf1,
 \qquad
 \mathcal N_0^{-1/2}\partial_k\Theta_{\mathrm{SR}}=\mathbf e_k.
\]
The central anticommutator of these four constant directions is the Minkowski
metric.  Every second spacetime derivative of $\Theta_{\mathrm{SR}}$ vanishes,
so it also solves the standard source-free constant-coefficient second-order
master operator in this flat branch.

More generally, for any constant Lorentz tetrad $\bar E_\mu$,
\[
 \Theta_{\mathrm{aff}}(x)=\Theta_0+
 \sqrt{\mathcal N_0}\,\bar E_\mu x^\mu
\]
generates that tetrad exactly.  This closes the special-relativistic
representer problem, although it does not prove uniqueness or curved-space
existence.

In polar coordinates, or in an accelerating frame,$\Gamma$ and $\Omega$ may
be nonzero even though spacetime is flat.  The invariant statement of flatness
is vanishing curvature, not vanishing connection coefficients in every frame.

\section{Four covariant gradients as a tetrad}
Define
\[
 E_\mu:=\frac{1}{\sqrt{\mathcal N_0}}D_\mu\Theta.
\]
The four objects $E_0,E_1,E_2,E_3$ are the local UBT rulers and clock direction.
They are not four extra matter fields; they are the first covariant derivatives
of the single fundamental field $\Theta$.

Let $i$ be the ordinary commuting complex unit and
$\mathbf e_1,\mathbf e_2,\mathbf e_3$ the quaternion units.  In the classical
Lorentz sector write
\[
 E_\mu=i e_\mu{}^0\mathbf1+e_\mu{}^k\mathbf e_k,
 \qquad e_\mu{}^a\in\mathbb R.
\]
Quaternion conjugation changes the signs of the three quaternion units but
leaves the complex scalar coefficient untouched:
\[
 E_\mu^\sharp=i e_\mu{}^0\mathbf1-e_\mu{}^k\mathbf e_k.
\]

\section{Why the anticommutator is a metric}
Quaternion multiplication contains a scalar product and a cross product.  In
the symmetrised product the cross products cancel:
\[
 \boxed{
 \frac12\left(E_\mu^\sharp E_\nu+E_\nu^\sharp E_\mu\right)
 =g_{\mu\nu}\mathbf1.}
\]
The right-hand side contains $\mathbf1$, the unit biquaternion.  The entire
left-hand side is already central: an ordinary real number multiplied by the
algebra unit.  No trace, real-part map, phase projection, or fiber average is
required.

Expanding the coefficients gives
\[
 g_{\mu\nu}
 =-e_\mu{}^0e_\nu{}^0+
  e_\mu{}^1e_\nu{}^1+
  e_\mu{}^2e_\nu{}^2+
  e_\mu{}^3e_\nu{}^3,
\]
which is exactly
\[
 g_{\mu\nu}=e_\mu{}^a e_\nu{}^b\eta_{ab},
 \qquad \eta_{ab}=\operatorname{diag}(-1,1,1,1).
\]
The minus sign of time comes from $i^2=-1$.

\section{What does the other half of the product do?}
The antisymmetric algebraic part is
\[
 \Sigma_{\mu\nu}
 =\frac12\left(E_\mu^\sharp E_\nu-E_\nu^\sharp E_\mu\right).
\]
It describes an oriented plane or bivector and is a natural carrier for local
rotation, boost, and spin information.  It is not itself the curvature.

Three expressions must be kept separate:
\begin{align*}
 \{E_\mu,E_\nu\}_\sharp
 &\longrightarrow g_{\mu\nu},\\
 [E_\mu,E_\nu]_\sharp
 &\longrightarrow \Sigma_{\mu\nu},\\
 [D_\mu,D_\nu]
 &\longrightarrow \mathcal R_{\mu\nu}.
\end{align*}

\section{Curvature from noncommuting derivatives}
For a one-sided connection, the curvature has the familiar form
\[
 \mathcal R_{\mu\nu}
 =\partial_\mu\Omega_\nu-\partial_\nu\Omega_\mu
  +[\Omega_\mu,\Omega_\nu].
\]
The final commutator is essential because biquaternion multiplication is
noncommutative.  Geometrically, curvature measures how a frame changes after
transport around a small closed loop.

The equation $E_\mu=D_\mu\Theta/\sqrt{\mathcal N_0}$ creates an additional
integrability condition.  This is the point at which the representation of the
connection becomes decisive.

\section{Why the old rank problem is disappearing locally}
A symmetric $4\times4$ metric has ten independent components.  Looking only at
one value of $\Theta\in\mathbb B\simeq\mathbb R^8$ suggests the misleading
comparison $8<10$.

The metric is not built from the value of $\Theta$ alone.  It is built from four
covariant derivatives $E_\mu$.  In the Lorentz sector these form a real
$4\times4$ tetrad $e_\mu{}^a$ with sixteen components.  Six changes are local
rotations and boosts that leave the metric unchanged.  Therefore
\[
 16-6=10.
\]
More rigorously, every small symmetric metric variation $h_{\mu\nu}$ is
produced by
\[
 \delta e_\mu{}^a=\frac12 h_{\mu\rho}e^{\rho a}.
\]
The tetrad-to-metric map therefore has rank ten at every nondegenerate tetrad.
This resolves the local kinematic rank question.  It does not yet prove that
every curved tetrad is generated by an on-shell UBT field.

\subsection{The right calculation can answer a poorly posed architectural question}
The former single-section rank obstruction was mathematically correct in an
induced-metric framework where Einstein equations were to be recovered only
from normal variations of one embedding-like section.  It did not constitute a
no-go theorem for the covariant tetrad $E_\mu=D_\mu\Theta/\sqrt{\mathcal N_0}$.
The compact-fiber construction repaired the former formulation by adding
profile directions, but the tetrad formulation shows that this repair was not
required by the original architecture.

This distinction motivates a general research rule:

\begin{quote}
Before adding fields, modes, dimensions, projections, averages, or auxiliary
structures to solve an obstruction, first test whether the obstruction is an
artefact of the formulation in which it was derived.
\end{quote}

The fiber route remains mathematically consistent in its own framework.  Its
problem is weak selection and representer redundancy, not an algebraic
contradiction.

\section{The problem of integrability}
Once the classical connection is reconstructed from the tetrad and torsion,
the defining equation becomes schematically
\[
 E_\mu=\mathcal N_0^{-1/2}
 \left(\partial_\mu\Theta+\rho_*(\Omega_\mu[E,T])\Theta\right).
\]
The unknown $E$ occurs on both sides.  This is an implicit nonlinear first-order
PDE or fixed-point system.  Eliminating $\Omega$ kinematically does not by
itself prove existence or uniqueness of a pair $(\Theta,E)$.

\subsection{One-sided obstacle leading to flatness}
Suppose the connection acts only from the left,
\[
 D^L_\mu\Theta=\partial_\mu\Theta+A_\mu\Theta,
\]
so that
\[
 [D^L_\mu,D^L_\nu]\Theta=F^L_{\mu\nu}\Theta.
\]
In a torsion-free compatible tetrad branch, antisymmetrising the covariant
Hessian gives
\[
 F^L_{\mu\nu}\Theta=0.
\]
If $\Theta$ is invertible as a biquaternion, then
\[
 F^L_{\mu\nu}=0.
\]
Thus the naive one-sided regular representation cannot generically support
curved GR while retaining an invertible $\Theta$ and torsion-free tetrad
compatibility.  Noninvertible stabilizer configurations may evade this result,
but they cannot be assumed to represent arbitrary geometries.

\subsection{Natural two-sided biquaternion derivative}
Biquaternions naturally admit both left and right multiplication.  Consider
\[
 \boxed{
 D_\mu\Theta=\partial_\mu\Theta+A_\mu\Theta-\Theta B_\mu.}
\]
A direct expansion gives
\[
 \boxed{
 [D_\mu,D_\nu]\Theta
 =F^A_{\mu\nu}\Theta-\Theta F^B_{\mu\nu}.}
\]
Torsion-free compatibility therefore requires
\[
 F^A_{\mu\nu}\Theta=\Theta F^B_{\mu\nu}.
\]
For invertible $\Theta$,
\[
 \boxed{
 F^A_{\mu\nu}=\Theta F^B_{\mu\nu}\Theta^{-1}.}
\]
Unlike the one-sided case, neither curvature must vanish.  The field $\Theta$
intertwines the left and right curvature representations.  This is a genuine
structural narrowing of GAP-10I: the two-sided/bimodule derivative is the
minimal algebra-native route that avoids the generic flatness obstruction.

\subsection{Reduction without a new field, torsion-free no-go and local right inversion with torsion}
On the Lorentz slice $W_L=\{X\mid X^\ddagger=-X\}$, preservation of the slice
and of its quadratic form reduces the pure pair, modulo a common central
one-form that cancels, to
\[
 \boxed{A_\mu=\Omega_\mu,\qquad B_\mu=-\Omega_\mu^\ddagger.}
\]
Thus $A_\mu$ and $B_\mu$ are not independent gravitational fields in this
branch.  They are the two module actions of the one spin connection.

With zero torsion, the same involution-compatible derivative makes a
nondegenerate solution $D_\mu\Theta=\sqrt{\mathcal N_0}E_\mu$ imply
\[
 \boxed{\mathring\nabla_\mu V^\nu=\delta_\mu{}^\nu,}
 \qquad \mathcal L_Vg=2g.
\]
The non-flat Schwarzschild vacuum exterior with nonzero mass has no such proper
homothety.  The pure pair is therefore closed kinematically and its $K=0$
branch is closed as a torsion-free no-go.

The qualifier matters.  On a sufficiently small non-null Gaussian patch,
choose $\rho\ne0$ with
$g^{\mu\nu}\partial_\mu\rho\partial_\nu\rho=\varepsilon=\pm1$ and set
\[
 V^\mu=\varepsilon\rho\nabla^\mu\rho,
 \qquad
 W_{\mu\nu}=g_{\mu\nu}-\mathring\nabla_\mu V_\nu,
\]
\[
 \boxed{K_{\nu\mu\rho}
 =\frac{W_{\mu\nu}V_\rho-V_\nu W_{\mu\rho}}{V^2}.}
\]
Then $K_{\nu\mu\rho}=-K_{\rho\mu\nu}$ and
$\nabla^{(\mathring\Gamma+K)}_\mu V^\nu=\delta_\mu{}^\nu$.  Therefore
\[
 \Theta=\sqrt{\mathcal N_0}V^a\mathbf u_a
\]
obeys $D_\mu\Theta=\sqrt{\mathcal N_0}E_\mu$.  Every smooth tetrad has such a
local representer with composite metric-compatible contortion and no
independent $A_\mu,B_\mu$ fields.

A relative term $P_\mu X-XQ_\mu$ is no longer required for local kinematic
existence.  It remains a possible torsion-free completion and, if used, must be
composite or auxiliary rather than a new propagating field.  A common central
constant does not remove the torsion-free obstruction because it cancels.

\section{Conditional dynamic partial closures}
\subsection{Algebraic Cartan torsion equation}
In the minimal first-order Hilbert--Palatini branch, independent variation of
the Lorentz connection gives
\[
 \epsilon_{abcd}T^a\wedge e^b=\kappa\tau_{cd}.
\]
For a nondegenerate tetrad this is a $24\times24$ pointwise linear map of rank
24. Thus zero spin current gives $T=0$, while a specified spin current gives a
unique torsion and contorsion. This closes the torsion algebra inside the
minimal branch, not the derivation of that branch from fundamental UBT.

\subsection{Propagation of the Lorentz section as a set of fixed points}
The anti-linear involution
\[
 \mathcal JX=-\overline{X^\sharp}
\]
has the physical Lorentz module as its fixed set. If a well-posed unique
Cauchy evolution is $\mathcal J$-equivariant and its data and sources are fixed
by $\mathcal J$, uniqueness implies that the solution remains in the Lorentz
slice. The final UBT action must still be checked against these hypotheses.

\subsection{Stability of the metric in imaginary time}
A vertical flow
\[
 \partial_\psi e_\mu{}^a=\Lambda_\psi{}^a{}_b e_\mu{}^b,
 \qquad \Lambda_{\psi ab}=-\Lambda_{\psi ba},
\]
is tangent to a local Lorentz gauge orbit and satisfies
$\partial_\psi g_{\mu\nu}=0$ identically. A second sufficient mechanism is a
unique $\psi$-translation-invariant evolution with $\psi$-independent initial
data. Selecting the physical mechanism and excluding unstable non-gauge modes
remains a dynamical task.

\subsection{Extended holonomy for prescribed curved coefficients}
For prescribed $E_\mu,A_\mu,B_\mu$, the inhomogeneous system
\[
 \partial_\mu\Theta+A_\mu\Theta-\Theta B_\mu
 =\sqrt{\mathcal N_0}E_\mu
\]
becomes homogeneous parallel transport for $Y=(\Theta,1)^T$ in an augmented
bundle. A single-valued solution with initial value $Y_0$ exists exactly when
the augmented holonomy fixes $Y_0$; flat augmented curvature is sufficient on
a simply connected domain. This closes prescribed-coefficient integrability,
not self-consistent UBT selection of the coefficients.

\subsection{Conditional Einstein-Lambda endpoint}
The minimal Palatini branch with zero spin current yields
\[
 G_{\mu\nu}+\Lambda g_{\mu\nu}=\kappa T_{\mu\nu}.
\]
Under the ordinary four-dimensional Lovelock assumptions, every local natural
divergence-free second-order metric equation without extra light geometric
fields has gravitational part $aG_{\mu\nu}+bg_{\mu\nu}$. Therefore the
conditional low-energy endpoint is unique. UBT must still derive those
assumptions, coefficients and matter terms from its canonical action.

\section{Implicit versus transcendent equations}
The words\emph{implicit} and\emph{transcendental} describe different
properties.

An equation is implicit when the unknown occurs inside the map that is supposed
to determine it:
\[
 E=\mathcal F(E,\Theta).
\]
An equation is transcendental when it contains nonalgebraic functions such as
an exponential, logarithm, trigonometric function, or Jacobi theta function:
\[
 E=e^{-E}.
\]
The UBT tetrad equation is already implicit and differential.  If the allowed
field $\Theta(q,\tau)$ is restricted to a Jacobi-theta or another transcendental
function class, then the complete system may be both implicit and
transcendental.  In a formal paper these properties should be stated
separately, even though intuitively both make the unknown ``entangled with
itself.''

\section{Derivatives versus variations}
A derivative such as $D_\mu\Theta$ describes how one field configuration changes
through spacetime.  A variation compares nearby possible configurations:
\[
 \Theta_\varepsilon=\Theta+\varepsilon\,\delta\Theta.
\]
The variation is
\[
 \delta\Theta=
 \left.\frac{d\Theta_\varepsilon}{d\varepsilon}\right|_{0}.
\]
Variations are needed when deriving field equations from an action.  They are
not replacements for spacetime derivatives.  When the connection depends on
the field,
\[
 \delta(D_\mu\Theta)
 =D_\mu(\delta\Theta)+\rho_*(\delta\Omega_\mu)\Theta
\]
and the induced connection variation must be included.

\section{What is closed and what remains open}
The following results are closed at their stated levels:
\begin{enumerate}
 \item the central anticommutator metric and rank-ten tetrad map;
 \item GAP-10$\Omega$-KIN/GR: reconstruction of the compatible connection and
       its Levi--Civita torsion-free branch;
 \item GAP-10T-PALATINI (conditional): the minimal Cartan torsion equation is
       algebraic and invertible;
 \item GAP-10L-CONN and GAP-10L-SYM (conditional): compatible transport and
       unique equivariant evolution preserve the Lorentz slice;
 \item GAP-10I-SR and GAP-10I-PRESCRIBED: the affine flat representer and the
       exact augmented-holonomy criterion for prescribed curved coefficients;
 \item GAP-10I-1S: the naive one-sided invertible curved route is a conditional
       no-go;
 \item GAP-10I-PAIR-KIN: the pure compatible pair reduces to the one spin
       connection;
 \item GAP-10I-PAIR-GR: the $K=0$ pure pair is a concurrent-vector no-go
       for the non-flat Schwarzschild vacuum exterior with nonzero mass;
 \item GAP-10I-TORSION-LOCAL: every smooth tetrad has a local single-$\Theta$
       representer with explicit composite metric-compatible contortion;
 \item GAP-10D-PALATINI/UNIQUENESS (conditional): the minimal first-order
       branch and Lovelock hypotheses select Einstein--$\Lambda$;
 \item GAP-10$\psi$-KIN/SYM: Lorentz-gauge flow or translation symmetry
       protects the classical metric under the stated hypotheses.
\end{enumerate}

The undivided full-theory gaps are narrowed rather than fully closed:
\begin{enumerate}
 \item GAP-10T-SPIN is closed conditionally for the direct fixed-background
       Palatini variation, while GAP-10T-FLAT-NOGO and
       GAP-10T-PAIRING-NOGO are closed as no-go results;
 \item GAP-10T-DYN: compute the full composite $\Theta$ variation and derive
       the selected first-order branch, a canonical torsion cancellation or
       translational/relative completion, and normalization from canonical UBT;
 \item GAP-10I-CURVED: local kinematics is closed; derive canonical
       action-level selection and physical admissibility of the composite
       torsionful branch, or an optional torsion-free relative bimodule branch,
       and establish regularity and global continuation of
       $(\Theta,E,A,B,T)$;
 \item GAP-10L-DYN: prove equivariance and well-posed uniqueness for the final
       complete dynamics;
 \item GAP-10D: derive the low-energy Palatini/Lovelock assumptions, couplings
       and matter action from the fundamental theory;
 \item GAP-10$\psi$: select the stability mechanism and exclude physical
       unstable non-gauge modes.
\end{enumerate}
GAP-B-MASTER and GAP-U2$\Theta$ remain open. The earlier compact-$\psi$ fiber
closure remains a mathematically consistent historical candidate completion,
but its weak selection and representer redundancy keep it outside the minimal
canonical route.
